\section{Infrastructure and Sensing Requirements}
\label{sec:hardware-requirements}

\subsection{Planning Principles}

Hardware planning supports \textbf{governed local processing}---not procurement of
surveillance capacity.
All selections must satisfy:

\begin{itemize}[noitemsep]
  \item ephemeral sensing with no default raw-media archive on DUALEXIS paths;
  \item no biometric processing (no facial recognition, voiceprints, or gait identification);
  \item operation on segmented school networks without cloud dependency for nominal awareness;
  \item physical security and asset accountability under school IT policies.
\end{itemize}

Indicative capacity tiers and vendor-neutral specifications for IT procurement appear in
Appendix~\ref{app:technical-architecture}.

\subsection{Capability Tiers (Institutional Summary)}

\begin{table}[htbp]
  \centering
  \caption{Capability tiers for pilot planning discussions (not binding SKUs).}
  \label{tab:capability-tiers}
  \small
  \begin{tabular}{@{}p{3.2cm}p{5.0cm}p{4.8cm}@{}}
    \toprule
    Tier & Typical zones & Planning note \\
    \midrule
    Light &
      Audio or sensor-only zones &
      Lower processing footprint; suitable for exit/door monitoring pilots. \\
    Standard &
      Single-stream video per zone &
      Supports zone-level activity descriptors; most MVP schools. \\
    Enhanced &
      Multi-stream or dense zones &
      Higher local processing; requires IT thermal/power planning. \\
    \bottomrule
  \end{tabular}
\end{table}

Schools should right-size tiers to \textbf{protocol mapping needs}, not maximum camera coverage.

\subsection{Camera and Sensing Policy}

\begin{itemize}[noitemsep]
  \item cameras cover \textbf{zones} (hallways, cafeterias, exit lobbies)---not individual
        identification targets;
  \item mounting avoids non-public areas unless legally justified and documented;
  \item transparency signage and stakeholder communications align with DPIA materials;
  \item existing school camera infrastructure may be reused only if governance approves
        scope and retention boundaries.
\end{itemize}

\textbf{Prohibited:} facial recognition modules, person re-identification, or persistent
video archives on DUALEXIS processing paths unless legally separated from this pilot scope.

\subsection{Audio and Environmental Inputs (Optional)}

\begin{itemize}[noitemsep]
  \item ambient audio may support \textbf{zone-level} stress indicators only;
  \item no speaker identification or speech-to-text of named individuals in MVP unless
        separately approved with strict purpose limitation;
  \item door, occupancy, and environmental sensors may supplement video where available.
\end{itemize}

\subsection{Site Coordination and Staff Access Host}

The site coordination service and staff interface require a \textbf{reliable on-premises
host} sized for event volume, review queues, and audit retention agreed in the DPIA.
MVP pilots may begin with a single coordinated host; extended pilots should plan
availability with school IT.

Requirements at planning level:
\begin{itemize}[noitemsep]
  \item encrypted storage for events, context records, and audit logs;
  \item backup procedure aligned with retention policy;
  \item UPS or equivalent for brief outages where feasible;
  \item locked facility access per school physical-security standards.
\end{itemize}

\subsection{Procurement Guidance for Leadership}

Steering committees should treat hardware as \textbf{enabling infrastructure} subordinate to:
\begin{enumerate}[noitemsep]
  \item governance sign-off and DPIA;
  \item staff training and protocol mapping;
  \item tabletop validation before live occupant sensing.
\end{enumerate}

Binding specifications for IT tender documents are in Appendix~\ref{app:technical-architecture}.
